Tras mirar el artículo te diría:

1.- Ser cuidadoso con no incurrir en autoplagio. Ahora mismo no recuerdo con exactitud el artículo de revista para saber si esto pasa.
2. Poner más bibliografía y las imágenes a dos columnas con lo que ya estamos en long paper.
3. Dar más importancia en un congreso como este a la parte de evaluación HCAI de la herramienta.


 Entiendo que para ahondar en la parte de evaluación y manifestarlo así desde el título del artículo sería interesante:
1.- En la introducción plantear que hay una carencia de metodologías/técnicas para evaluar los sistemas HGML
 2.- En el Related Work aclarar definitivamente esa carencia ya que lo hecho en la comunidad científica no lo aborda
3.- Establecer una nueva sección donde se haga una cierta aportación diciendo cómo se propone evaluar los sistemas HGML y ya sigues en la siguiente sección presentando tu sistema y aplicando la evaluación (que tendría algo de novedoso).
 
En definitiva, tratar de evitar que el revisor tenga la impresión de que la aportación ya estaba plenamente hecha con los criterios HCAI y simplemente se han aplicado. Durante esta extensión del deadline (hay 10 días). Ver si se puede escribir una pequeña sección proponiendo algo que mejore/amplíe lo ya establecido en cuanto a criterios HCAI. Esto sería un contenido directamente llevado a tu tesis.